\documentclass[10pt]{beamer}

\usepackage{verbatim}
%\usepackage[latin1]{inputenc}
\usepackage{amssymb,amsmath,amsfonts,url}
\usepackage{graphicx,color}
\usepackage{pgf,pgfarrows,pgfnodes,pgfautomata,pgfheaps,pgfshade}
\usepackage[utf8]{inputenc}
\usepackage[french]{babel}
\usepackage{tikz}
\usetikzlibrary{positioning}
\usepackage[all]{xy}
\usepackage{amstext}
\usepackage{xspace}
\usepackage{mdwtab}


\mode<article> {
  \usepackage{fullpage}
  \usepackage{hyperref}
  \setjobnamebeamerversion{advancedcourse.beamer}
}





\newcommand{\alertonce}[1]{\addtocounter{beamerpauses}{-1}\alert<+>{#1}}
\newcommand{\alerttwice}[1]{\addtocounter{beamerpauses}{-1}\alert<+-+(1)>{#1}}

%\newcommand{\subsubsubsection}[1]{\noindent\textbf{#1}}

\newcommand{\cfbox}[1]{\begin{center}\fbox{#1}\end{center}}
\newcommand{\cbox}[1]{\begin{center} #1 \end{center}}

%%%%%%%%%%%%%%%%%%%%%%%%%%%%%%%%%%%%%%%%%%%%%%%%%%%%%%%%%%%%%%%%%%%%%%%%%%%%%%%%%%%%%%%%%

\graphicspath{{Graphics/}}
%\input def_0.tex

%\input{StyleSlides.tex}
\mode<presentation> {
    \usepackage{beamerthemeBoadilla}

  %\beamertemplatetransparentcovereddynamic
  \beamertemplateballitem
}



\title[IN100 Cours 9]{\textbf{Fondements informatiques I}\\
Cours 9: Algorithmes : piles et files}

\author[]{Sorina Ionica \texttt{sorina.ionica@uvsq.fr} \\~\\Sandrine Vial  \texttt{sandrine.vial@uvsq.fr} }


\institute{}

\date{}


%\setbeamercolor{normal text}{fg=black}

\begin{document}


\frame{\titlepage}



\begin{frame}[fragile]{File}

Une file (en anglais queue) est une structure de données basée sur le principe que les premiers éléments ajoutés dans la file seront les premiers à en être retirés.


\vspace{0.5 cm}


\begin{tikzpicture}[
    box/.style={draw, minimum width=0.6cm, minimum height=0.6cm, anchor=west},
    node distance=0cm,
    >=stealth
]

% === Initial queue ===
\node[box] (a1) at (0,0) {3};
\node[box, right=of a1] (a2) {5};
\node[box, right=of a2] (a3) {7};

%\node[draw=none, anchor=west] at (0,-0.8) {Initial queue};

% === Queue after enqueue 4 and 6 (same level) ===
\node[box] (b1) at (5,0) {\textcolor{red}{6}};
\node[box, right=of b1] (b2) {\textcolor{red}{4}};
\node[box, right=of b2] (b3) {3};
\node[box, right=of b3] (b4) {5};
\node[box, right=of b4] (b5) {7};


% Arrow from initial queue to enqueue queue
\node[draw=none, anchor=west] at (2.5,0.2) {\small{Rajouter 4, 6}};
\draw[->, thick] (a3.east) -- (b1.west);

% === Queue after two dequeue operations ===
\node[box] (c1) at (6.3,-2) {6};
\node[box, right=of c1] (c2) {4};
\node[box, right=of c2] (c3) {3};
\node[box, right=of c3] (c4) {5};

\node[draw=none, anchor=west] at (8.8,-1.8) {\small{Enlever}};


\node[box] (d1) at (10.3,-2) {6};
\node[box, right=of d1] (d2) {4};
\node[box, right=of d2] (d3) {3};

\draw[->, thick] (c4.east) -- (d1.west);

%\node[draw=none, anchor=west] at (6,-3.8) {After dequeue twice};

% Arrow from enqueue queue to dequeue queue
\node[draw=none, anchor=west] at (6.5,-0.7) {\small{Enlever}};
\draw[->, thick] (b5.east) .. controls +(0.9,-2.2) and +(-0.8,2.2) .. (c1.west);

\end{tikzpicture}

\pause

\vspace{0.3 cm}

Nombreuses applications: 
\begin{itemize}
\item traiter dans l'ordre une liste de transactions, la gestion d'un stock
\item gestion des tâches dans l'ordre d'arrivée par exemple sur un système d'exploitation 
\item transmission des  données dans l'ordre sur un réseau
\end{itemize}


\end{frame}


\begin{frame}[fragile]{File}


\vspace{0.5 cm}

\begin{tikzpicture}[
    box/.style={draw, minimum width=0.6cm, minimum height=0.6cm, anchor=west},
    node distance=0cm,
    >=stealth
]

% === Initial queue ===
\node[box] (a1) at (0,0) {3};
\node[box, right=of a1] (a2) {5};
\node[box, right=of a2] (a3) {7};

%\node[draw=none, anchor=west] at (0,-0.8) {Initial queue};

% === Queue after enqueue 4 and 6 (same level) ===
\node[box] (b1) at (5,0) {\textcolor{red}{6}};
\node[box, right=of b1] (b2) {\textcolor{red}{4}};
\node[box, right=of b2] (b3) {3};
\node[box, right=of b3] (b4) {5};
\node[box, right=of b4] (b5) {7};


% Arrow from initial queue to enqueue queue
\node[draw=none, anchor=west] at (2.5,0.2) {\small{Rajouter 4, 6}};
\draw[->, thick] (a3.east) -- (b1.west);

% === Queue after two dequeue operations ===
\node[box] (c1) at (6,-2) {6};
\node[box, right=of c1] (c2) {4};
\node[box, right=of c2] (c3) {3};
\node[box, right=of c3] (c4) {5};

\node[draw=none, anchor=west] at (8.5,-1.8) {\small{Enlever}};


\node[box] (d1) at (10,-2) {6};
\node[box, right=of d1] (d2) {4};
\node[box, right=of d2] (d3) {3};

\draw[->, thick] (c4.east) -- (d1.west);

%\node[draw=none, anchor=west] at (6,-3.8) {After dequeue twice};

% Arrow from enqueue queue to dequeue queue
\node[draw=none, anchor=west] at (6.5,-0.7) {\small{Enlever}};

\draw[->, thick] (b5.east) .. controls +(0.9,-2.2) and +(-0.8,2.2) .. (c1.west);

\end{tikzpicture}


\vspace{0.3cm}

Pour mettre en place une file nous avons besoin de deux procédures (fonctions en Python): 
\begin{itemize}
\item une pour ajouter en tête de file 
\item une pour enlever de la file. 
\end{itemize}

\vspace{-0.25cm}
\pause

\begin{columns}
\column{5 cm}
\begin{exampleblock}{}
\begin{verbatim}
file=[3,5,7]
ajouter(file,4)
ajouter(file,6)
enlever(file)
enlever(file)
\end{verbatim}
\end{exampleblock}
\end{columns}


\end{frame}



\frame[containsverbatim]{
\frametitle{Ajouter dans le file}

On suppose que notre pile est représentée par un tableau dynamique. 

\medskip
\begin{alertblock}{Procédure ajouter}
\begin{verbatim}
Input: La file f et l'élément e à ajouter
Output : La file modifiée 

l<-longueur(f)
On agrandit f d'une case (c.a.d. on alloue de la mémoire, 
longueur(f)<- longueur(f)+1)
Pour tout i allant de l-1 à 0
     f[i+1]<- f[i]
     
f[0]=e
Return f
\end{verbatim}
\end{alertblock}


}


\frame[containsverbatim]{
\frametitle{Enlever de la file}

On suppose que notre pile est représentée par un tableau dynamique. 

\medskip
\begin{alertblock}{Procédure enlever}
\begin{verbatim}
Input: La file f 
Output : Le dernier élément de la file et  la file modifiée 

l<-longueur(f)

Si (l>0) alors: 
       e<-f[l-1]
       On rétrécit f d'une case (désallouer mémoire, 
       longueur(f)<-l-1)
Return e,f
\end{verbatim}
\end{alertblock}


}


\begin{frame}[fragile]
{Une file en Python}

En Python, on utilisera les listes pour représenter des tableaux dynamiques. 

\begin{exampleblock}{}
\begin{columns}
\column{5 cm}
\begin{verbatim}
def ajouter(f,e):
    l=len(f)
    f.append(0)
    for i in range(l):
   	   f[l-i]=f[l-i-1]
    f[0]=e
\end{verbatim}

\pause

\column{5 cm}
\vspace{-1.5 cm}
\begin{verbatim}
def enlever(f):
    if (len(f)>0):
        return f.pop()
\end{verbatim}
\end{columns}
\end{exampleblock}


\pause

\vspace{0.5 cm}


\vspace{0.2 cm}

La fonction \texttt{pop(i)} enlève l'élément d'indice \texttt{i} dans la liste. \\

\vspace{0.2 cm}

L'appel \texttt{pop()} enlève le dernier élément de la liste. 

\vspace{0.2 cm}
Serait-il possible d'utiliser la fonction \texttt{remove} au lieu de \texttt{pop}?  \\


\end{frame}


\begin{frame}[fragile]
{Une file en Python}



\begin{exampleblock}{}
\begin{columns}
\column{5 cm}
\vspace{-0.275 cm}
\begin{verbatim}
def ajouter(f,e):
    l=len(f)
    f.append(0)
    for i in range(l):
   	   f[l-i]=f[l-i-1]
    f[0]=e
\end{verbatim}

\column{5 cm}
\vspace{-1.5 cm}
\begin{verbatim}
def enlever(f):
    if (len(f)>0):
       return f.pop()
\end{verbatim}
\end{columns}
\end{exampleblock}

\vspace{0.5 cm}

Pour quoi ces fonctions ne renvoient pas \texttt{f} comme demandé? 
\pause

\begin{exampleblock}{}
\begin{verbatim}
f=[3,5,7]
ajouter(f,4)
ajouter(f,6)    
print("voici la file",f)
enlever(f)
print("voici la file",f)
\end{verbatim}
\end{exampleblock}



\end{frame}





\begin{frame}[fragile]{Exemple d'utilisation}

\'Ecrire un programme qui permet de  traiter les  transactions sur un compte bancaire pendant une journée. 
%Les recettes seront ajoutées au montant du compte, les dépenses seront déduites sur le compte. 
A la fin de la journée, le montant restant sur le compte est affiché. 

%Input: Un tableau tab qui contient la liste des transactions de la journée et le montant initial sur le compte 
%Output: Le montant restant sur le compte en fin de journée 

\vspace{0.5 cm}
En Python, on utilisera  la fonction \texttt{enlever}. 


\begin{exampleblock}{}
\begin{verbatim} 
montant=3498
f=[53,-78,198,-23,-243]

while len(f)>0: 
   transaction=enlever(f)
   montant=montant+transaction
   
print("montant restant", montant)    
\end{verbatim}
\end{exampleblock}


\end{frame}


\begin{frame}[fragile]{Pile}

\vspace{-0.5 cm}

Une pile (en anglais stack) est une structure de données  fondée sur le principe "dernier arrivé, premier sorti"  \\
Le  dernier élément ajouté à la pile est le premier à en sortir. 
  
  \vspace{0.8 cm}
  
  
  \begin{tikzpicture}[
  rect/.style={draw, minimum width=0.5cm, minimum height=0.3cm, align=center},
  stack/.style={inner sep=0pt},
  >=stealth
]

% === First stack (Initial) at x=0 ===
\begin{scope}[xshift=0cm]
  % stack from bottom (y=0) to top (y=2*0.8)
  \node[rect] (s1_1) at (0,0) {3};
  \node[rect] (s1_2) at (0,0.4) {5};
  \node[rect] (s1_3) at (0,0.8) {7};
\end{scope}

% === Second stack (after pushes) at x=5cm ===
\begin{scope}[xshift=2.5cm]
  \node[rect] (s2_1) at (0,0) {3};
  \node[rect] (s2_2) at (0,0.4) {5};
  \node[rect] (s2_3) at (0,0.8) {7};
  \node[rect] (s2_4) at (0,1.2) {\textcolor{red}{4}}; % pushed 4
  %\node[rect] (s2_5) at (0,1.6) {6}; % pushed 6 (top)
\end{scope}

% === Third stack (after two pops) at x=10cm ===
\begin{scope}[xshift=5cm]
  \node[rect] (s3_1) at (0,0) {3};
  \node[rect] (s3_2) at (0,0.4) {5};
  \node[rect] (s3_3) at (0,0.8) {7}; % back to original top
   \node[rect] (s3_4) at (0,1.2) {4}; % pushed 4
  \node[rect] (s3_5) at (0,1.6) {\textcolor{red}{6}}; % pushed 6 (top)
\end{scope}

\begin{scope}[xshift=7.5cm]
  \node[rect] (s4_1) at (0,0) {3};
  \node[rect] (s4_2) at (0,0.4) {5};
  \node[rect] (s4_3) at (0,0.8) {7}; % back to original topdépiler
   \node[rect] (s4_4) at (0,1.2) {4}; % pushed 4
\end{scope}

\begin{scope}[xshift=10cm]
  \node[rect] (s5_1) at (0,0) {3};
  \node[rect] (s5_2) at (0,0.4) {5};
  \node[rect] (s5_3) at (0,0.8) {7}; % back to original top

\end{scope}


% === Arrows between stacks ===
% arrow from top of initial stack to top of after-push stack
\draw[->, thick] (s1_3.east) .. controls +(right:8mm) and +(left:8mm) .. (s2_4.west)
  node[midway,above] {\small{empiler 4}};

% arrow from after-push stack to after-pop stack
\draw[->, thick] (s2_4.east) .. controls +(right:8mm) and +(left:8mm) .. (s3_5.west)
  node[midway,above] {\small{empiler 6}};
  
  
\draw[->, thick] (s3_5.east) .. controls +(right:8mm) and +(left:8mm) .. (s4_4.west)
  node[midway,above] {\small{dépiler}};

\draw[->, thick] (s4_4.east) .. controls +(right:8mm) and +(left:8mm) .. (s5_3.west)
  node[midway,above] {\small{dépiler}};    

% optional: small braces to show stack boundaries (visual aid)
%\draw[dashed] (-0.9,-0.15) rectangle (0.9,2.05);
%\draw[dashed] (4.1,-0.15) rectangle (5.9,3.45);
%\draw[dashed] (9.1,-0.15) rectangle (10.9,2.05);

\end{tikzpicture}

  \vspace{0.8 cm}
  
  \pause


Nombreuses applications: 
\begin{itemize}
\item inverser une séquence 
\item évaluer des expressions arithmétiques (à voir plus tard)
\item dans les navigateurs internet pour sauvegarder les pages visitées 
\end{itemize}
 
 
\end{frame}





\frame[containsverbatim]{
\frametitle{Représenter une pile, empiler}

Par un tableau dynamique : sommet de pile = dernière case
\begin{itemize}
\item empiler = ajouter en dernière position
\item dépiler = supprimer l’élément en dernière position
\end{itemize}


\begin{alertblock}{Algorithme pour empiler}
\begin{verbatim} 
Input: Pile p et élément e
Output:  Pile p modifiée 

On agrandit p d'une case (allocation mémoire etc.) 
On rajoute e. 
Renvoyer p
\end{verbatim}
\end{alertblock}

}

\frame[containsverbatim]{
\frametitle{Dépiler}


\begin{alertblock}{Algorithme pour dépiler}
\begin{verbatim} 
Input: Pile p 
Output:  Pile p modifiée et la valeur dépilée

l<-longueur(f)
Si (l>0) alors: 
       e<-f[l-1]
On rétrécit p d'une case (désallouer mémoire etc.)

Renvoyer p,e
\end{verbatim}
\end{alertblock}



}



\frame[containsverbatim]{
\frametitle{Une pile en Python}

On utilisera une liste pour représenter la pile. 


\begin{exampleblock}{}
\begin{columns}
\column{5 cm}
\vspace{-0.55 cm}
\begin{verbatim} 
def empiler(p,e):
     p.append(e)
\end{verbatim}

\column{5 cm}

\vspace{-0.2 cm}
\begin{verbatim} 
def depiler(p):
   if len(p)>0:
       return p.pop()
\end{verbatim}
\end{columns}
\end{exampleblock}

\begin{exampleblock}{}

\begin{verbatim} 
p=[3,5,7]
empiler(p,4)
empiler(p,6)
print("voici l'état de la pile", p) #p=[3,5,7,4,6]
depiler(p)
depiler(p)
print("voici l'état de la pile", p) #p=[3,5,7]
\end{verbatim}

\end{exampleblock}

}





\begin{frame}[fragile]{Exemple de programme utilisant une pile}


Inverser un tableau ou une chaîne de caractères en utilisant une pile. \\

En Python, on utilise les fonctions \texttt{empiler} et \texttt{depiler}. 

\vspace{0.5 cm}

\begin{exampleblock}{En Python}
\begin{verbatim} 
a = [3, 5, 7, 2]
inverse = []

while (len(a)>0):
        empiler(inverse, depiler(a))

print(inverse) 
\end{verbatim}
\end{exampleblock}

\vspace{0.5 cm}


\end{frame}










\end{document}
